\subsection{Experimental Setup}
\label{subsec:experimental-setup}

We evaluate Neural Sheaf Diffusion (NSD) on three families of tasks designed
to test long-range information propagation: synthetic barbell graphs,
graph-transfer benchmarks, and large urban road networks. All experiments use
the implementation from \texttt{sheaf-mpnn}, pinned to a fixed Git revision,
and are executed through a common training and evaluation pipeline. We compare
four choices of restriction maps: general, orthogonal, diagonal, and identity.
The identity variant is an additional control in which every node--edge
restriction map is fixed to the identity matrix and is therefore not learned.
For each map family, we consider stalk dimensions
$d \in \{2,3,5\}$ and hidden dimensions
$h \in \{16,32\}$, giving $24$ model configurations. Unless stated otherwise,
each configuration is evaluated using three random seeds.

\paragraph{Barbell regression.}
Following the synthetic over-squashing experiment of
Bamberger et al., each graph consists of two cliques connected by a sparse
set of bridge edges. Both the number of vertices per clique and the number of
bridges are configurable; the default graph contains ten vertices per clique
and one bridge. Each vertex has a $40$-dimensional input feature. Features in
the two cliques are sampled independently from
$\mathcal{U}(-\sqrt{3},0)$ and $\mathcal{U}(0,\sqrt{3})$, respectively. The
target at every vertex is the mean input feature of the opposite clique,
making successful prediction dependent on information passing through the
bottleneck. We use three NSD layers, Adam with learning rate $10^{-3}$, and
train for at most $500$ epochs. Performance is measured using the mean,
over vertices, of the squared Euclidean prediction error.

\paragraph{Graph-transfer tasks.}
We implement the Ring, CrossedRing, CliquePath, and binary Tree tasks from
Di Giovanni et al. Each graph contains a designated source and target vertex.
The target receives one of five one-hot class vectors, the source receives the
zero vector, and all remaining vertices receive a constant vector. The model
must recover the target class at the source. Ring, CrossedRing, and CliquePath
are evaluated at graph sizes
$N \in \{2,6,10,20,30\}$, while binary trees are evaluated at depths
$r \in \{2,\ldots,8\}$. The number of NSD layers is set to the shortest-path
distance between source and target, preventing under-reaching while retaining
the effect of over-squashing. For each topology and size, we generate balanced
sets of $5000$, $500$, and $500$ graphs for training, validation, and testing.
Models are trained with Adam, learning rate $10^{-3}$, weight decay
$5 \times 10^{-4}$, and batch size $128$. We report source-node
classification accuracy and macro-F1. A separate legacy configuration is
retained solely for comparison with the original implementation; all primary
results use validation-based model selection and a held-out test set.

\paragraph{City networks and reproducibility.}
Finally, we evaluate NSD on the Paris and Shanghai datasets from
City-Networks. Their road graphs contain $114{,}127$ and $183{,}917$ vertices,
respectively, with $37$ augmented input features and ten classes derived from
local eccentricity. We retain the canonical $10\%/10\%/80\%$
train/validation/test masks and use $16$ NSD layers to match the long-range
nature of the labels. Training uses AdamW with learning rate
$5 \times 10^{-4}$, weight decay $10^{-5}$, and dropout $0.2$; checkpoints
are selected by validation accuracy. Experiments are run on cloud GPUs using
mixed bfloat16 precision when supported. Every run stores its resolved
configuration, random seed, software and Git revisions, epoch-level metrics,
best checkpoint, and final test results. Synthetic datasets are generated
deterministically, while city data are cached in a persistent volume. This
ensures that all reported results can be reconstructed from the recorded run
identifier and configuration.

\paragraph{Information-flow and geometric diagnostics.}
For a trained model \(Z=f_\theta(X,G)\), we quantify the influence of source
vertex \(u\) on target vertex \(v\) using the Jacobian block
\(J_{v,u}=\partial z_v/\partial x_u\). We report its entrywise
\(\ell_1\) norm, Frobenius norm, spectral norm, and the gradient norm of the
ground-truth logit. These quantities are aggregated into total and mean
influence at each shortest-path distance, with normalized decay curves,
bootstrap confidence intervals, decay slopes, and an influence-weighted
radius. Synthetic graphs use exact full blocks; for the city graphs,
ground-truth-logit vector--Jacobian products are evaluated on a deterministic
sample of test vertices, while full matrix and singular-value diagnostics are
retained for a smaller subset. We additionally decompose each end-to-end
Jacobian into products of local layer Jacobians along monotone shortest-path
computational walks. This separates the sum of geodesic contributions, a
canonical-path contribution, and the remaining non-geodesic component, while
also measuring cancellation between paths. For NSD, we record the singular
spectra and condition numbers of restriction and normalized transport maps at
every layer. Finally, we compare the Ollivier--Ricci curvature of the original
unweighted graph with that of the learned effective metric
\(\ell_e=(\epsilon+\omega_e)^{-1}\), using the
\texttt{OTDSinkhornMix} approximation with \(\alpha=0.5\). All diagnostics are
stored alongside both initial and trained checkpoints, permitting direct
analysis of how learning changes information flow, anisotropy, bottleneck
strength, curvature, and predictive performance.
